\begin{notechapter}
\chapter{Convolution and Image Filtering}\label{ch:convolution-filtering}
\chaptersubtitle{Linear shift-equivariant maps, and the filters of classical vision}

\begin{chapterintro}
Convolution is the operation that convolutional networks are built on, and long before
them it was the basic tool of signal and image processing. This chapter defines it
precisely, distinguishes it from cross-correlation (which is what deep-learning libraries
actually compute), proves the property that makes it the right operation for images, and
uses it to compute image gradients, the starting point of classical feature extraction. You
will learn
\begin{itemize}[nosep]
  \item convolution and cross-correlation in one and two dimensions, their properties, and
        their matrix form;
  \item why convolutions are exactly the linear maps that commute with translations;
  \item boundary handling, separable kernels, smoothing and derivative filters;
  \item what the Sobel kernel computes, with the right factor and sign;
  \item gradient magnitude and direction, the Harris corner detector and HOG features.
\end{itemize}
\prerequisites{linear algebra (matrices, eigenvalues of symmetric matrices); partial
derivatives.}
\end{chapterintro}

\section{Convolution and cross-correlation}\label{convolution-filtering:sec:definitions}

Signals are sequences $x:\Z\to\R$, zero outside a finite support.

\begin{definition}[Convolution and cross-correlation]\label{convolution-filtering:def:conv}
The \emph{convolution} and the \emph{cross-correlation} of $x$ with a kernel $k$ are
\[
  (x * k)[n] = \sum_mx[m]\,k[n - m],\qquad\qquad(x\star k)[n] = \sum_mx[m + n]\,k[m] .
\]
\end{definition}

Convolution flips the kernel before sliding it; cross-correlation slides it unflipped. With
the flipped kernel $\tilde k[m] = k[-m]$, $x\star k = x * \tilde k$: the two operations differ only by the orientation of
the kernel.

\begin{proposition}[Algebraic properties]\label{convolution-filtering:prop:algebra}
Convolution is bilinear, commutative ($x * k = k * x$) and associative
($(x * k) * h = x * (k * h)$), with the unit impulse $\delta$ ($\delta[0] = 1$, $0$ elsewhere) as neutral element.
Cross-correlation is bilinear but not commutative: $(k\star x)[n] = (x\star k)[-n]$.
\end{proposition}

\begin{proof}
Commutativity: substitute $m' = n - m$. Associativity: both sides equal
$\sum x[a]\,k[b]\,h[c]$ over the triples with $a + b + c = n$. For cross-correlation,
\[
  (k\star x)[n] = \sum_mk[m + n]\,x[m] = \sum_{m'}k[m']\,x[m' - n] = (x\star k)[-n] .
\]
\end{proof}

\begin{example}[A short signal]\label{convolution-filtering:ex:short}
Let $x = (x_1, x_2, x_3)$ at positions $1,2,3$ and $k = (k_1, k_2)$ at positions $1,2$.
\[
  \begin{array}{c|c}
    n & (x * k)[n] = (k * x)[n]\\\hline
    2 & x_1k_1\\
    3 & x_1k_2 + x_2k_1\\
    4 & x_2k_2 + x_3k_1\\
    5 & x_3k_2
  \end{array}
  \qquad\qquad
  \begin{array}{c|c|c}
    n & (x\star k)[n] & (k\star x)[n]\\\hline
    -2 & 0 & k_1x_3\\
    -1 & x_1k_2 & k_1x_2 + k_2x_3\\
    0 & x_1k_1 + x_2k_2 & k_1x_1 + k_2x_2\\
    1 & x_2k_1 + x_3k_2 & k_2x_1\\
    2 & x_3k_1 & 0
  \end{array}
\]
The two correlation columns are mirror images of each other. In the convolution, $k_1$
always multiplies the later sample: the kernel has been flipped.
\end{example}

\paragraph{Matrix form.} For fixed $k$, $x\mapsto x * k$ is linear, so on finite signals it is a matrix.
With $x\in\R^3$ and $k = (k_1, k_2)$, the ``full'' convolution, of length $3 + 2 - 1 = 4$, is
\[
  \begin{bmatrix}(x * k)[2]\\(x * k)[3]\\(x * k)[4]\\(x * k)[5]\end{bmatrix}
  = \begin{bmatrix}k_1 & 0 & 0\\ k_2 & k_1 & 0\\ 0 & k_2 & k_1\\ 0 & 0 & k_2\end{bmatrix}
  \begin{bmatrix}x_1\\x_2\\x_3\end{bmatrix} .
\]
The matrix is \emph{Toeplitz}: constant along its diagonals. Software offers three output
sizes for a signal of length $n$ and a kernel of length $k$: \emph{full} ($n + k - 1$ outputs, every
overlap), \emph{valid} ($n - k + 1$, only full overlaps), and \emph{same} ($n$, centred).

\section{Shift equivariance}\label{convolution-filtering:sec:equivariance}

Let $(\tau_ax)[n] = x[n - a]$ be the translation by $a$. A map $T$ is \emph{shift-equivariant} if $T(\tau_ax) = \tau_a(Tx)$
for all $x$ and $a$: translating the input translates the output.

\begin{theorem}[Convolutions are the linear shift-equivariant maps]\label{convolution-filtering:thm:lti}
For every kernel $k$, $x\mapsto x * k$ is linear and shift-equivariant. Conversely, every linear
shift-equivariant map $T$ on finitely supported signals is a convolution: $Tx = x * h$ with
$h = T\delta$, its \emph{impulse response}.
\end{theorem}

\begin{proof}
$(\tau_ax * k)[n] = \sum_mx[m - a]k[n - m] = \sum_{m'}x[m']k[n - a - m'] = (x * k)[n - a]$. Conversely, write
$x = \sum_mx[m]\,\tau_m\delta$. By linearity and equivariance, $Tx = \sum_mx[m]\,\tau_m(T\delta) = \sum_mx[m]\,\tau_mh$, whose value at $n$ is
$\sum_mx[m]h[n - m] = (x * h)[n]$.
\end{proof}

This is the reason for convolutions in vision: an edge, an eye or a wheel should be
detected in the same way wherever it appears. A layer that is linear and treats all
positions alike \emph{must} be a convolution. The same holds for cross-correlation, which is a
convolution with a flipped kernel.

\begin{pitfall}[Convolution in deep learning]
The ``convolution'' layers of PyTorch, TensorFlow and JAX compute a cross-correlation.
Since the kernel is learned, this is harmless: learning $k$ or its flip is the same problem.
It matters when a hand-designed kernel such as Sobel is applied with library functions:
\texttt{scipy.signal.convolve2d} flips it, \texttt{scipy.signal.correlate2d},
\texttt{cv2.filter2D} and \texttt{torch.nn.functional.conv2d} do not, and the results differ by a sign for
antisymmetric kernels (\cref{convolution-filtering:ex:sobel}).
\end{pitfall}

\section{Images}\label{convolution-filtering:sec:images}

A grey-level image is a function $I[i,j]$ of the row $i$ (increasing downwards) and the column $j$
(increasing to the right). Two-dimensional cross-correlation with a kernel $K$ indexed by
offsets $u,v\in\{-r,\dots,r\}$ is
\[
  (I\star K)[i,j] = \sum_{u=-r}^r\sum_{v=-r}^rI[i + u, j + v]\,K[u,v],
\]
the sum of the products of the kernel with the $(2r + 1)\times(2r + 1)$ patch centred at $(i,j)$. Convolution uses
$K[-u,-v]$, the kernel rotated by $180^\circ$.

\paragraph{Borders.} Near the border the patch leaves the image. Options are: compute only
where it fits (``valid'', the output shrinks by $2r$ in each dimension); pad with zeros (the
default in networks); replicate the border pixels; reflect the image; or wrap it around
(circular). Zero padding creates artificial edges at the border; reflection does not.

\paragraph{Separable kernels.} If $K = uv^\top$ is an outer product, $I\star K = (I\star u)\star v^\top$: a column filter followed
by a row filter. The cost per pixel drops from $(2r + 1)^2$ to $2(2r + 1)$ multiplications.

\paragraph{Smoothing.} The box filter $\frac19\mathbf 1\mathbf 1^\top$ averages a $3\times3$ patch. The binomial filter
\[
  \tfrac1{16}\begin{bmatrix}1\\2\\1\end{bmatrix}\begin{bmatrix}1&2&1\end{bmatrix} = \tfrac1{16}\begin{bmatrix}1&2&1\\2&4&2\\1&2&1\end{bmatrix}
\]
approximates a Gaussian; Gaussian filters
$K[u,v]\propto e^{-(u^2 + v^2)/(2s^2)}$ are separable and isotropic, and are the standard way to remove noise before
differentiating.

\section{Image gradients}\label{convolution-filtering:sec:gradients}

Think of the image as samples of a smooth intensity $I(x,y)$, with $x$ the column and $y$ the row
coordinate. The central difference $\frac12(I[i,j + 1] - I[i,j - 1])$ approximates $\partial I/\partial x$ with error $O(h^2)$; it is the
cross-correlation with $\frac12\begin{bmatrix}-1 & 0 & 1\end{bmatrix}$. Differences amplify noise, so they are combined with
smoothing in the other direction.

\begin{definition}[Sobel, Prewitt and Scharr kernels]\label{convolution-filtering:def:sobel}
\[
  S_x = \begin{bmatrix}1\\2\\1\end{bmatrix}\begin{bmatrix}-1 & 0 & 1\end{bmatrix} = \begin{bmatrix}-1 & 0 & 1\\-2 & 0 & 2\\-1 & 0 & 1\end{bmatrix},\qquad
  S_y = S_x^\top = \begin{bmatrix}-1 & -2 & -1\\0 & 0 & 0\\1 & 2 & 1\end{bmatrix} .
\]
Prewitt uses the smoothing weights $(1,1,1)$, Scharr $(3,10,3)$, for better rotational accuracy.
\end{definition}

\begin{proposition}[What Sobel computes]\label{convolution-filtering:prop:sobel}
On a linear image $I[i,j] = aj + bi + c$, the cross-correlations are $I\star S_x = 8a$ and $I\star S_y = 8b$ at every interior
pixel: Sobel by cross-correlation computes $8\,\partial I/\partial x$ and $8\,\partial I/\partial y$, with $y$ pointing down. (Prewitt gives
$6\times$, Scharr $32\times$ the derivative.) Convolution with $S_x$ gives $-8\,\partial I/\partial x$.
\end{proposition}

\begin{proof}
Each row of $S_x$ computes $I[i + u,j + 1] - I[i + u, j - 1] = 2a$; the weights $1,2,1$ sum to $4$. Likewise for $S_y$.
Convolution flips $S_x$, which changes its sign.
\end{proof}

The factor is irrelevant for detecting edges (thresholds are relative) but matters when
the gradient is used quantitatively. From the two responses $G_x = I\star S_x$, $G_y = I\star S_y$:
\[
  G = \sqrt{G_x^2 + G_y^2},\qquad\theta = \operatorname{atan2}(G_y, G_x)\in(-180^\circ, 180^\circ] .
\]
$G$ is large on edges; $\theta$ is the direction of steepest increase, perpendicular to the edge.
The two-argument arctangent distinguishes opposite directions, which $\arctan(G_y/G_x)$ cannot,
and handles $G_x = 0$. For display, $\theta$ is mapped to the hue of the HSV colour wheel (red $0^\circ$, green
$120^\circ$, blue $240^\circ$), so that opposite directions get complementary colours.

\begin{example}[An edge with a step]\label{convolution-filtering:ex:sobel}
The $5\times5$ image, dark ($0$) on the left and bright ($10$) on the right, with the boundary
moving one column left in the bottom two rows:
\[
  I = \begin{bmatrix}0&0&0&10&10\\0&0&0&10&10\\0&0&0&10&10\\0&0&10&10&10\\0&0&10&10&10\end{bmatrix},\qquad
  G_x = \begin{bmatrix}0&40&40\\10&40&30\\30&40&10\end{bmatrix},\qquad
  G_y = \begin{bmatrix}0&0&0\\10&20&10\\10&20&10\end{bmatrix}
\]
(valid cross-correlation, so the outputs are the $3\times3$ interior pixels). The vertical part
of the edge gives $G_x = 40 = 8\times(10/2)$, a jump of $10$ over two pixels; where the boundary steps, $G_y$
is positive (brighter below). The magnitudes and directions are
\[
  G = \begin{bmatrix}0&40&40\\14.1&44.7&31.6\\31.6&44.7&14.1\end{bmatrix},\qquad
  \theta = \begin{bmatrix}0^\circ&0^\circ&0^\circ\\45^\circ&26.6^\circ&18.4^\circ\\18.4^\circ&26.6^\circ&45^\circ\end{bmatrix}:
\]
the gradient points right, and tilts downwards near the step. \texttt{scipy.signal.convolve2d}
with $S_x$ returns $-G_x$.
\end{example}

\paragraph{Second derivatives.} The Laplacian kernel $\begin{bmatrix}0&1&0\\1&-4&1\\0&1&0\end{bmatrix}$ approximates
$\partial^2I/\partial x^2 + \partial^2I/\partial y^2$; edges are its zero crossings. The Canny detector (Canny, \emph{A
Computational Approach to Edge Detection}, IEEE TPAMI, 1986) smooths with a Gaussian,
computes $G$ and $\theta$, keeps local maxima of $G$ along $\theta$, and links them with two thresholds.

\section{From gradients to features}\label{convolution-filtering:sec:features}

\paragraph{Corners.} On an edge the gradient has one direction; at a corner, several. The
\emph{structure tensor} summarizes the gradients in a window $\mathcal W$ around a pixel, with weights
$w$ (often Gaussian):
\[
  M = \sum_{(i,j)\in\mathcal W}w_{ij}\begin{bmatrix}G_x^2 & G_xG_y\\ G_xG_y & G_y^2\end{bmatrix} .
\]
$M$ is symmetric positive semidefinite. Its eigenvalues $\lambda_1\geq\lambda_2\geq0$ measure the variation of the
image along its eigenvectors: both small in a flat region, $\lambda_1\gg\lambda_2\approx0$ on an edge, both large at
a corner. Harris and Stephens (\emph{A Combined Corner and Edge Detector}, Alvey Vision
Conference, 1988) avoid the eigenvalues with the response $R = \det M - \kappa(\tr M)^2 = \lambda_1\lambda_2 - \kappa(\lambda_1 + \lambda_2)^2$,
$\kappa\approx0.05$, which is large and positive at corners and negative on edges.

\paragraph{Histograms of oriented gradients.} HOG (Dalal and Triggs, CVPR, 2005) divides
the image into cells of, say, $8\times8$ pixels, builds in each cell a histogram of the directions $\theta$
weighted by the magnitudes $G$, normalizes the histograms over overlapping blocks of cells,
and concatenates them. The descriptor is robust to small shifts and to lighting changes;
with a linear SVM (\cref{ch:svm}) it was the standard pedestrian detector before CNNs.
SIFT (Lowe, \emph{Distinctive Image Features from Scale-Invariant Keypoints}, IJCV, 2004)
detects keypoints as extrema of differences of Gaussians across scales and describes each
by gradient histograms in a rotated local frame, which makes it invariant to scale and
rotation.

\paragraph{Learned filters.} The first layer of a convolutional network trained on natural
images learns, among others, oriented edge detectors similar to Sobel kernels, and
colour-blob detectors (Krizhevsky, Sutskever and Hinton, NeurIPS, 2012, Fig.~3). The
following chapters replace hand-designed filters by learned ones, stacked in many layers.

\paragraph{In code.}
\begin{lstlisting}[language=Python]
import numpy as np
from scipy.signal import correlate2d

Sx = np.array([[-1, 0, 1], [-2, 0, 2], [-1, 0, 1]], float)
Sy = Sx.T
Gx = correlate2d(I, Sx, mode="same", boundary="symm")   # 8 dI/dx
Gy = correlate2d(I, Sy, mode="same", boundary="symm")   # 8 dI/dy, y down
G = np.hypot(Gx, Gy)
theta = np.degrees(np.arctan2(Gy, Gx))                    # in (-180, 180]
# OpenCV: cv2.Sobel(I, cv2.CV_64F, 1, 0, ksize=3) gives the same Gx
\end{lstlisting}

\begin{chaptersummary}
\begin{itemize}[nosep]
  \item Convolution $\sum_mx[m]k[n - m]$ flips the kernel and is commutative and associative; cross-correlation
        $\sum_mx[m + n]k[m]$ does not flip, is not commutative, and is what deep-learning ``convolutions'' compute.
  \item Convolutions are exactly the linear shift-equivariant maps, determined by their impulse
        response; on finite signals they are Toeplitz matrices.
  \item Sobel combines a central difference with $(1,2,1)$ smoothing; by cross-correlation it gives $8\,\partial I/\partial x$
        and $8\,\partial I/\partial y$ ($y$ down), by convolution the opposite sign.
  \item Gradient magnitude and $\operatorname{atan2}$ direction feed edge detection (Canny), corner detection (the
        structure tensor and Harris response) and descriptors (HOG, SIFT).
\end{itemize}
\end{chaptersummary}

\section*{Exercises}
\addcontentsline{toc}{section}{Exercises}

\begin{exercise}\label{convolution-filtering:exr:1}
Compute $x * k$ and $x\star k$ for $x = (1, 2, 3)$ (positions $0,1,2$) and $k = (1, -1)$ (positions $0,1$). Which one computes
the backward difference $x[n] - x[n - 1]$, and what does the other compute?
\solutionlink{convolution-filtering:exr:1}
\end{exercise}

\begin{exercise}\label{convolution-filtering:exr:2}
Show that the transpose of the $4\times3$ Toeplitz matrix of $x\mapsto x * k$ (full convolution, $k = (k_1,k_2)$) is the
matrix of a valid cross-correlation with $k$, mapping $\R^4\to\R^3$.
\solutionlink{convolution-filtering:exr:2}
\end{exercise}

\begin{exercise}\label{convolution-filtering:exr:3}
Show that the Scharr kernel gives $32\,\partial I/\partial x$ on a linear image, and that the Laplacian kernel gives $0$ on
any image of the form $aj + bi + c$ and $2$ on $I[i,j] = j^2$.
\solutionlink{convolution-filtering:exr:3}
\end{exercise}

\begin{exercise}\label{convolution-filtering:exr:4}
Compute the structure tensor (all weights $1$) of the $3\times3$ block of $(G_x, G_y)$ in
\cref{convolution-filtering:ex:sobel}, its eigenvalues, and the Harris response with $\kappa = 0.05$. Then do the
same for the top row alone. Which looks more like a corner?
\solutionlink{convolution-filtering:exr:4}
\end{exercise}

\begin{exercise}\label{convolution-filtering:exr:5}
How many multiplications per pixel does a $7\times7$ Gaussian filter need directly, and as two
one-dimensional passes?
\solutionlink{convolution-filtering:exr:5}
\end{exercise}
\end{notechapter}
